\section*{Bab \ref{ch:b3:cytoskeleton-motility} --- Dinamika Sitoskeleton dan Gerak Sel}

\begin{solution}{exo:b3:cytoskeleton-motility:1}
Aktin: \qty{7}{nm}, aktin globular, polar (berduri/runcing), ATP;
korteks, penonjolan, kontraksi bersama \omterm{def:b3:cytoskeleton-motility:motors}{miosin}. \omterm{def:b3:cytoskeleton-motility:filaments}{Mikrotubulus}: \qty{25}{nm},
dimer $\alpha\beta$-tubulin, polar (plus/minus), GTP; angkutan
jarak jauh, gelendong, silia. \omterm{def:b3:cytoskeleton-motility:filaments}{Filamen antara}: \qty{10}{nm},
protein kumparan lingkar (keratin, vimentin, lamin), apolar, tanpa nukleotida;
kekuatan mekanis sel dan nukleusnya.
\end{solution}

\begin{solution}{exo:b3:cytoskeleton-motility:2}
Kepekatan monomer ketika sebuah ujung tidak tumbuh maupun menyusut, yakni
$k_{\text{lepas}}/k_{\text{ikat}}$. Kedua ujungnya berbeda karena subunitnya
menghidrolisis ATP-nya sesudah penyisipan, sehingga kedua ujungnya menampilkan
spesi kimia yang berbeda (aktin-ATP yang tiba di ujung berdurinya, aktin-ADP
yang pergi dari ujung runcingnya) dengan kegemaran yang berbeda. Dengan
\omterm{thm:b3:cytoskeleton-motility:treadmilling}{kepekatan kritis} yang sama akan ada satu kesetimbangan tunggal: tanpa \omterm{thm:b3:cytoskeleton-motility:treadmilling}{gerak ban
berjalan}, tanpa fluks, dan tanpa cara membayar pergantian yang terarah.
\end{solution}

\begin{solution}{exo:b3:cytoskeleton-motility:3}
Ke tepinya pada \omterm{def:b3:cytoskeleton-motility:filaments}{mikrotubulus}: \omterm{def:b3:cytoskeleton-motility:motors}{kinesin}, menuju ujung plusnya. Ke pusatnya:
\omterm{def:b3:cytoskeleton-motility:motors}{dinein} sitoplasma, menuju ujung minusnya. Serabut tegangan: \omterm{def:b3:cytoskeleton-motility:motors}{miosin} II,
menuju ujung berduri aktinnya. \omterm{def:b3:cytoskeleton-motility:cilia}{Silium}: \omterm{def:b3:cytoskeleton-motility:cilia}{dinein aksonem}, menuju
ujung minus ganda tetangganya (yakni pangkalnya).
\end{solution}

\begin{solution}{exo:b3:cytoskeleton-motility:4}
Penonjolan (Rac bagi \omterm{def:b3:cytoskeleton-motility:migration}{lamelipodium}, Cdc42 bagi filopodium), adhesi
(\omterm{def:b3:cytoskeleton-motility:migration}{integrin}; di hilir Rac dan Rho), kontraksi (Rho, lewat
\omterm{def:b3:cytoskeleton-motility:motors}{miosin} II), penarikan bagian belakangnya (pelepasan adhesinya, yang digerakkan
kontraksinya).
\end{solution}

\begin{solution}{exo:b3:cytoskeleton-motility:5}
$C_{c}^{+} = 1/5 = \qty{0.2}{\micro M}$, $C_{c}^{-} = 0.8/1 =
\qty{0.8}{\micro M}$; $C_{\text{ss}} = (1 + 0.8)/(5 + 1) = \qty{0.3}{\micro
M}$; $J = 5\times 0.3 - 1 = 0.5$ subunit tiap detik, \qty{1.35}{nm/s}.
\end{solution}

\begin{solution}{exo:b3:cytoskeleton-motility:6}
$800/8 = 100$ ATP tiap detik. Satu meter memakan $1/(8\times 10^{-7}) =
1.25\times 10^{6}$ s, kira-kira dua minggu, dan $1.25\times 10^{8}$ langkah dan
ATP. Difusi akan memakan \num{16000} tahun: selisih semiliar kali lipat.
\end{solution}

\begin{solution}{exo:b3:cytoskeleton-motility:7}
$F_{\max} = 8.3\times 10^{-20}/3.6\times 10^{-8} = \qty{2.3}{pN}$. Energi yang
sama yang tersebar sepanjang langkah lebih panjang memberi gaya lebih kecil ---
sebuah tuas yang lebih panjang. \omterm{def:b3:cytoskeleton-motility:motors}{Kinesin} (langkah pendek, \qty{6}{pN}) cocok
bagi muatan berat; \omterm{def:b3:cytoskeleton-motility:motors}{miosin} V (langkah panjang) menempuh jarak lebih jauh tiap
ATP dan cocok bagi angkutan cepat yang ringan.
\end{solution}

\begin{solution}{exo:b3:cytoskeleton-motility:8}
(a) Taksol membekukan \omterm{def:b3:cytoskeleton-motility:filaments}{mikrotubulusnya}: gelendongnya tak dapat menelusuri,
menangkap, dan membetulkan, sehingga mitosisnya terhenti pada titik periksanya;
sel yang merayap terus menonjol tetapi kehilangan kekutubannya untuk jangka
panjang. (b) Kolkisina membuang \omterm{def:b3:cytoskeleton-motility:filaments}{mikrotubulusnya}: tanpa gelendong, mitosisnya
terhenti; perayapannya berlanjut pada aktinnya tetapi tanpa arah yang tetap.
(c) Sitokalasin menudungi ujung berdurinya: tanpa penonjolan, tanpa perayapan;
mitosisnya berlanjut tetapi sitokinesisnya gagal, sehingga memberi sel
berinti dua. (d) Penghambatan Rac: tanpa \omterm{def:b3:cytoskeleton-motility:migration}{lamelipodium}, sehingga tanpa
perayapan; pembelahannya sebagian besar tak terpengaruh.
\end{solution}

\begin{solution}{exo:b3:cytoskeleton-motility:9}
Waktu tinggalnya seorde dengan waktu paruhnya, \qty{20}{s} (tetapan waktunya
$20/\ln 2 \approx \qty{29}{s}$); bagian yang tak bergerak \qty{10}{\%}.
Polimerisasi pada tepinya harus memasok kemajuannya sekaligus aliran
mundurnya: $1 + 1.5 = \qty{2.5}{\micro m/min}$.
\end{solution}

\begin{solution}{exo:b3:cytoskeleton-motility:10}
Ujung yang tumbuh menahan sebuah tudung tubulin-GTP; di belakangnya hidrolisis
menghasilkan tubulin-GDP yang tegang. Hilangnya tudung itu --- peristiwa acak
yang peluangnya naik seiring melambatnya pertumbuhan --- melepaskan tegangannya
dan ujungnya menyusut cepat sampai tudung baru terbentuk. Tepat di atas
\omterm{thm:b3:cytoskeleton-motility:treadmilling}{kepekatan kritisnya}, tiap \omterm{def:b3:cytoskeleton-motility:filaments}{mikrotubulus} pada tiap saat entah bertudung dan
tumbuh entah tanpa tudung dan runtuh, sehingga panjangnya menyimpang menjadi
populasi yang panjang dan populasi yang melenyap, bukan mengerumun pada sebuah
rata-rata. Selnya memperoleh penjelajahan volumenya yang cepat serta kemampuan
memantapkan secara terpilih --- sebuah ujung plus yang mencapai kinetokor atau
tapak korteks ditangkap lalu dipertahankan, sedangkan yang lain didaur ulang
dalam hitungan menit: telusur dan tangkap.
\end{solution}

\begin{solution}{exo:b3:cytoskeleton-motility:11}
Pada skala ini inersianya dapat diabaikan dan persamaan fluidanya
berbalikan-waktu: sebuah gerak yang menapaki kembali jalurnya (sebuah dayung
yang keluar lalu kembali) mengembalikan tubuhnya ke tempat semula, seberapa
cepat pun tiap sapuannya. Sebuah heliks yang berputar tak pernah menapaki
kembali jalurnya --- geraknya kiral dan menerus --- sehingga ia menghasilkan
dorongan neto. Pada \qty{100}{Hz} dengan seribu proton tiap putaran, jumlahnya
$10^{5}$ proton tiap detik tiap motornya.
\end{solution}

\begin{solution}{exo:b3:cytoskeleton-motility:12}
Silia saluran napas tak dapat berdenyut, lendir dan bakterinya tak
dibersihkan, dan infeksinya berulang. Flagelum sperma, yang dibangun pada
\omterm{def:b3:cytoskeleton-motility:cilia}{aksonem} yang sama, tak bergerak: kemandulan. Di dalam embrio, silia bergerak
pada nodusnya menggerakkan aliran ke kiri yang menetapkan sumbu kiri--kanannya;
tanpa aliran itu sisinya dipilih secara acak, sehingga separuh pasiennya
berorgan terbalik.
\end{solution}

\begin{solution}{pb:b3:cytoskeleton-motility:1}
\textbf{1.} $\qty{300}{\micro m^3} = 3\times 10^{-13}$ L; $\times
2\times 10^{-4}$ mol/L $= 6\times 10^{-17}$ mol $= 3.6\times 10^{7}$
molekul; $1.8\times 10^{7}$ di dalam filamen.
\textbf{2.} $1.8\times 10^{7}\times \qty{2.7}{nm} = \qty{4.9}{cm}$
filamen; pada \qty{0.25}{\micro m} masing-masing, sekitar $2\times 10^{5}$
filamen.
\textbf{3.} Protein penudung menghalangi sebagian besar ujung berdurinya, dan
profilin serta timosin-$\beta$4 mengikat monomernya sehingga aktin-ATP yang
benar-benar bebas berada dekat \omterm{thm:b3:cytoskeleton-motility:treadmilling}{kepekatan kritisnya}; pertumbuhannya terkurung
pada ujung yang dibuka tudungnya oleh selnya.
\textbf{4.} $C_{\text{ss}} = 2.2/12.9 = \qty{0.17}{\micro M}$; $J = 11.6
\times 0.17 - 1.4 = 0.6$ subunit tiap detik; \qty{1.6}{nm/s} $=
\qty{0.1}{\micro m/min}$.
\textbf{5.} $1/0.1 = 10$ kali lipat.
\textbf{6.} \qty{1}{\micro m/min} $= \qty{16.7}{nm/s} = 6.2$ subunit tiap
detik tiap filamen; $\times 2\times 10^{5} = 1.2\times 10^{6}$ ATP tiap
detik, kira-kira $0.1\,\%$ pergantian selnya.
\textbf{7.} $1/(8\times 10^{-7}) = 1.25\times 10^{6}$ s $\approx 14.5$
hari; $1.25\times 10^{8}$ langkah dan ATP.
\textbf{8.} $L^{2}/2D$: $(1)^{2}/(2\times 10^{-12}) = 5\times 10^{11}$ s
$\approx \num{16000}$ tahun bagi satu meter; $(10^{-5})^{2}/(2\times
10^{-12}) = 50$ s bagi \qty{10}{\micro m}. Difusi melayani di dalam sebuah
badan sel; apa pun yang lebih panjang menuntut motor.
\textbf{9.} $5\times 10^{4}/6.02\times 10^{23} = 8.3\times 10^{-20}$ J
$= 20\,k_{B}T$.
\textbf{10.} $8.3\times 10^{-20}/8\times 10^{-9} = \qty{10}{pN}$;
daya gunanya $6/10 \approx 60\,\%$.
\textbf{11.} $F = 6\pi\times 0.01\times 10^{-7}\times 8\times 10^{-7} =
1.5\times 10^{-14}$ N $= \qty{0.015}{pN}$, empat ratus kali di bawah
\omterm{prop:b3:cytoskeleton-motility:physics}{gaya hentinya}: satu motor sudah berlimpah melawan seretan kental; muatan
yang sungguhan adalah rintangan dan tambatan.
\textbf{12.} $10^{4}\times 100 = 10^{6}$ ATP tiap detik, $0.1\,\%$ anggaran
neuronnya: angkutannya murah. Namun motornya berhenti begitu ATP-nya
turun, sehingga kegagalan mitokondria melaparkan sinapsnya akan segala yang
dibuat di dalam badan selnya --- yakni kegagalan angkutan akson yang terlihat
pada penyakit kemunduran saraf.
\textbf{13.} \qty{10}{\micro m/min} $= \qty{167}{nm/s}$; $167/2.7 = 62$
subunit tiap detik tiap filamen.
\textbf{14.} $200\times 20 = 4000$ filamen; $4000\times 62 = 2.5\times
10^{5}$ subunit tiap detik.
\textbf{15.} $2.5\times 10^{5}$ ATP tiap detik, $0.025\,\%$ anggarannya.
\textbf{16.} $\qty{3}{\micro m}/(\qty{10}{\micro m/min}) = \qty{0.3}{min}
= \qty{18}{s}$; $2.5\times 10^{5}\times 18 = 4.5\times 10^{6}$ subunit
di dalam jalanya.
\textbf{17.} $1\,\text{pN}\times 2.7\,\text{nm} = \qty{2.7}{pN.nm} =
0.66\,k_{B}T$: riak sebesar energi itu terjadi dengan peluang
kira-kira $e^{-0.66} \approx 0.5$ --- celah sebesar satu subunit membuka
terus-menerus, dan ratcetnya sepenuhnya masuk akal.
\textbf{18.} Polimerisasinya sendiri yang mendorong, dengan bakterinya
menjadi ``membran'' di depan jalanya sendiri; dengan permesinan Arp2/3 selnya,
ia mencapai laju \omterm{def:b3:cytoskeleton-motility:migration}{lamelipodium} dan lebih --- sampai
\qty{0.5}{\micro m/s}, yakni \qty{30}{\micro m/min}.
\textbf{19.} Pada $C = K_{d}$, $\theta = 1/2$ dan $\mathrm{d}\theta/
\mathrm{d}C = 1/(4C)$: selisih \qty{2}{\%} pada $C$ memberi $\Delta\theta
= 0.005$; dengan \num{25000} reseptor tiap sisi, $125$ lagi yang terisi di
bagian depannya.
\textbf{20.} Tiap sisinya memiliki sekitar \num{12500} yang terisi, yang
beriak sebesar $\sqrt{\num{12500}} \approx 110$, sehingga selisih kedua
sisinya beriak sekitar $160$: angka $125$ tenggelam di dalam deraunya pada
saat mana pun. Reseptornya mengikat ulang kira-kira sekali tiap detik;
perata-rataan sepanjang $N$ detik mengecilkan deraunya sebesar $\sqrt{N}$ ---
sepuluh detik membawanya ke $50$ dan landaiannya pun terbaca.
\textbf{21.} Rac (bersama Cdc42) di depan, Rho di belakang. Dengan
Rac aktif di mana-mana, selnya menonjol ke segala arah, melebar, dan tak ke mana-mana.
\textbf{22.} $30/3 = \qty{10}{\micro m/min}$: seperti yang diberikan.
\textbf{23.} Penonjolannya digantikan penggelembungan yang digerakkan tekanan;
adhesi dan kontraksinya masih bergantung pada aktin dan \omterm{def:b3:cytoskeleton-motility:motors}{miosin}, sehingga aktin
tetap mutlak bagi sisa daurnya.
\textbf{24.} Bagian depannya dibangun ulang tiap \qty{18}{s}: pergeseran
keaktifan Rac mengarahkan ulang polimerisasinya dalam satu umur jala, lalu
bagian depan yang baru memimpin, dan langkah lambat di belakangnya menyusul.
\textbf{25.} \omterm{thm:b3:cytoskeleton-motility:treadmilling}{Gerak ban berjalan} di luar sel \qty{1.6}{nm/s}; sekitar $14.5$ hari
untuk menempuh aksonnya; sekitar $2.5\times 10^{5}$ ATP tiap detik bagi
penonjolannya.
\end{solution}
